\subsection{Proof of the Chern-Weil Theorem}

We will now proof Theorem \ref{chern_weil:thm:chern_weil_thm}. The hard part of the proof is to show that $\eta(\Omega)$ is closed. For this we basically do the computation outlined in \cite{emilia_globale_ana} but add the missing details and fix a few mistakes. We need the following two Lemmas.

\begin{lemma}
    \label{chern_weil:lem:technical_lemma1}
    Let $\overline{\eta} \in I_r(k)$ and $X_1,\ldots,X_r,Y \in \R^{k \times k}$. Then 
    \[
    \sum_{i=1}^{r} \overline{\eta}(X_1,\ldots,[X_i,Y],\ldots,X_r) = 0,
    \]
    where $[X_i,Y] = X_iY - YX_i$ is the commutator.
\end{lemma}

\begin{proof}
    Let $A(t) \in \mathrm{GL}(k)$ be a smooth curve with $A(0) = \mathrm{id}$ and $\dot{A}(0) = Y$. Then 
    \begin{align*}
         \frac{d}{dt}|_{t=0} A(t)^{-1}X_iA(t) &= \frac{d}{dt}|_{t=0}A(t)^{-1}X_i + \frac{d}{dt}|_{t=0}X_iA(t) \\ &= -YX_i + X_iY \\ &= [X_i,Y],
    \end{align*}
    where we used that 
    \[
    0 = \frac{d}{dt}|_{t=0} A(t)A(t)^{-1} = \frac{d}{dt}|_{t=0} A(t) + \frac{d}{dt}|_{t=0} A(t)^{-1} = Y + \frac{d}{dt}|_{t=0} A(t)^{-1},
    \]
    such that 
    \[
    \frac{d}{dt}|_{t=0} A(t)^{-1} = -Y.
    \]
    The invariance of $\overline{\eta}$ now yields 
    \begin{align*}
        0 &= \frac{d}{dt}|_{t=0} \, \overline{\eta}(A(t)^{-1}X_1A(t),\ldots,A(t)^{-1}X_rA(t)) \\
        &= \sum_{i=1}^{r}\overline{\eta}(X_1,\ldots,\frac{d}{dt}|_{t=0}A(t)^{-1}X_iA(t),\ldots,X_r)\\
        &= \sum_{i=1}^{r} \overline{\eta}(X_1,\ldots,[X_i,Y],\ldots,X_r).
    \end{align*}
\end{proof}

Before stating the second Lemma let us define the commutator of two differential forms $\Omega_1 \in \Omega^{l_1}(M,\R^{k \times k})$ and $\Omega_2 \in \Omega^{l_2}(M,\R^{k \times k})$ by 
\[
[\Omega_1,\Omega_2] := \Omega_1 \wedge \Omega_2 - (-1)^{l_1l_2} \, \Omega_2 \wedge \Omega_1.
\]
If $\Omega_i = \omega_i \otimes A_i$ for some $\omega_i \in \Omega^{l_i}(M)$ and $A_i \in \Omega^0(M,\R^{k \times k})$ then this can also be written as 
\[
[\Omega_1,\Omega_2] = \omega_1 \wedge \omega_2 \otimes [A_1,A_2],
\]
such that this definition agrees with the one in \cite{emilia_globale_ana}, but again has the benefit of being coordinate independent.

\begin{lemma}[Bianchi Identity]
    \label{chern_weil:lem:bianchi_identity}
    Let $E \to M$ be a vector bundle with a connection, let $\Theta$ be the associated connection form (with respect to some local frame) and let $\Omega$ be the curvature form. Then 
    \[
    d \Omega + [\Theta,\Omega] = 0.
    \]
\end{lemma}
\begin{proof}
    This is just a matter of writing out the definitions.
\end{proof}

We are now ready to proof the Chern-Weil Theorem.

\begin{proof}[Proof of Theorem \ref{chern_weil:thm:chern_weil_thm}]
    Since closedness is a local propery we can directly use the local definition of $\eta(\Omega)$, i.e. we can assume that $\Omega$ is a matrix valued $2$-form. Let $\overline{\eta} = \Phi_r^{-1}(\eta) \in S_r(k)$. In order to compute $d \eta(\Omega) = d\overline{\eta}(\Omega,\ldots,\Omega)$, let us more generally show that 
    \begin{equation}
        \label{chern_weil:eq:equation1}
    d \overline{\eta}(\Omega_1,\ldots,\Omega_r) = \sum_{i=1}^{r} \overline{\eta}(\Omega_1,\ldots,d\Omega_i,\ldots,\Omega_r)
    \end{equation}
    for matrix valued $2$-forms $\Omega_i \in \Omega^2(M,\R^{k \times k})$. By linearity we may assume that all except for one entry in the matrix $\Omega_i = (\Omega_i^{jl})_{j,l=1}^k$ of $2$-forms vanish, so in particular
    \[
    \Omega_i = \omega_i \otimes E_i,
    \]
    where $E_i$ is a constant matrix (with only one non-zero entry being equal to $1$). We then compute 
    \begin{align*}
        d \overline{\eta}(\Omega_1,\ldots,\Omega_r) &= d \left(\sum_{i=1}^{r} \omega_1 \wedge \ldots \wedge \omega_r \overline{\eta}(E_1,\ldots,E_r)\right) \\
        &= \sum_{i=1}^{r} \omega_1 \wedge \ldots \wedge d\omega_i \wedge \ldots \omega_r \overline{\eta}(E_1,\ldots,E_r) \\
        &= \sum_{i=1}^{r} \overline{\eta}(\omega_1 \otimes E_1,\ldots, d\omega_i \otimes E_i, \ldots, \omega_r \otimes E_r) \\
        &= \sum_{i=1}^{r} \overline{\eta}(\Omega_1,\ldots,d\Omega_i,\ldots,\Omega_r)
    \end{align*}
    Note that $d E_i = 0, d \overline{\eta}(E_1,\ldots,E_r) = 0$ since both are constant. Next we claim that 
    \begin{equation}
        \label{chern_weil:eq:equation2}
    \sum_{i=1}^{r} \overline{\eta}(\Omega_1,\ldots,[\Theta,\Omega_i],\ldots, \Omega_r) = 0,
    \end{equation}
    where $\Theta$ is the connection form (or any matrix valued $1$-form for that matter). Again we may assume by linearity that $\Theta = \theta \otimes E$ for some constant matrix $E$, such that 
    \begin{align*}
        \overline{\eta}(\Omega_1,\ldots,[\Theta,\Omega_i],\ldots, \Omega_r) & = \overline{\eta}(\omega_1 \otimes E_1,\ldots, \theta \wedge \omega_i \otimes [E,E_i], \ldots, \omega_r \otimes E_r) \\
        & = \theta \wedge \omega_1 \wedge \ldots \wedge \omega_r \, \overline{\eta}(E_1,\ldots,[E,E_i],\ldots,E_r),
    \end{align*}
    where we used commutativity of the wedge product for $2$-forms. Taking the sum over $i$ the claim follows from Lemma \ref{chern_weil:lem:technical_lemma1}. Now we plug in $\Omega = \Omega_i$ for the curvature form $\Omega$ and compute 
    \begin{align*}
        d \eta(\Omega) & = d\overline{\eta}(\Omega,\ldots,\Omega) \\
            & \oset[4pt]{(\ref{chern_weil:eq:equation1})}{=} \sum_{i=1}^{r} \overline{\eta}(\Omega,\ldots,d\Omega,\ldots,\Omega) \\
            & \oset[4pt]{(\ref{chern_weil:eq:equation2})}{=} \sum_{i=1}^{r} \overline{\eta}(\Omega,\ldots,d\Omega + [\Theta,\Omega],\ldots,\Omega) = 0,
    \end{align*}
    where we used the Bianchi identity (Lemma \ref{chern_weil:lem:bianchi_identity}) in the last equation.

    Next we show the independence of the chosen connection. This is basically taken from \cite{rahman2017} but with a few more details added. Let $\nabla^0, \nabla^1$ be two different connections on $E$ with curvature forms $\Omega_0, \Omega_1$ and consider the pullback bundle $\pi^\ast E$ under the projection $\pi \colon M \times \R$. We want to define a connection $\nabla$ on $\pi^\ast E$ that smoothly interpolates between the connections $\nabla^0$ and $\nabla^1$. To this end let $\lambda \colon M \times \R \to \R$ be the projection and set 
    \[
    \nabla := \lambda \pi^\ast \nabla^1 + (1 - \lambda) \pi^\ast \nabla^0.
    \]
    It is easily seen that a convex combination of connections is again a connection (even if the coefficients are smooth functions). Now let $\iota_t \colon M \to M \times \R, \; p \mapsto (p,t)$ be the inclusions. We claim that $\iota_0^\ast \nabla = \nabla^0$ and $\iota_1^\ast \nabla = \nabla^1$. For this we have to check the characteristic property of pullback connections (cf. \ref{riem_mnf:prop:pullback_connections}). So let $s \in \Gamma(\pi^\ast E)$ and $v \in T_p M$. The claim now is that
    \[
    \nabla^0_v (\iota_0^\ast s) = \nabla_{(\iota_0)_\ast v} s.
    \]
    Let us abbreviate $w:= (\iota_0)_\ast v$. For the right hand side we have
    \[
    \nabla_{w} s = (\pi^\ast \nabla^0)_{w} s,
    \]
    where we used that $w$ is a tangent vector at $(p,0)$ and $\lambda(p,0) = 0$. We would like to replace $s$ with the section $\pi^\ast\iota_0^\ast s$, since then the claim just follows from the definition of the pullback connection: 
    \[
    (\pi^\ast \nabla^0)_{w} (\pi^\ast\iota_0^\ast s) = \nabla^0_v \iota_0^\ast s.
    \]
    In order to do this choose a local frame $e_1,\ldots,e_k$ of $E$ and write $s = s^i \pi^\ast e_i$ for smooth functions $s^i \colon M \times \R \to \R$. Then $\pi^\ast\iota_0^\ast s = (s^i \circ \iota_0 \circ \pi) \pi^\ast e_i$ and thus 
    \begin{align*}
    (\pi^\ast \nabla^0)_{w} (\pi^\ast\iota_0^\ast s) & = w \cdot (s^i \circ \iota_0 \circ \pi)  e_i(p) + s^i(p,0) \nabla^0_v e_i \\
        &= w \cdot s^i  e_i(p) + s^i(p,0) \nabla^0_v e_i \\
        &= (\pi^\ast \nabla^0)_{w} s.
    \end{align*}
    The same argument holds for $\iota_1$ and $\nabla^1$. In fact $\iota_t^\ast \nabla = \nabla^t$, where $\nabla^t = t\nabla^1 + (1-t)\nabla^0$. Now let $\Omega$ be the curvature form of $\nabla$, then by Lemma \ref{chern_weil:lem:naturality_of_curvature_form}
    \[
    \iota_0^\ast \Omega = \Omega_0, \; \iota_1^\ast \Omega = \Omega_1.
    \]
    Thus $\iota_0^\ast \eta(\Omega) = \eta(\Omega_0)$ and $\iota_1^\ast \eta(\Omega) = \eta(\Omega_1)$. Since the maps $\iota_0$ and $\iota_1$ are homotopic they induce the same map on cohomology, such that the statement follows.
    \end{proof}

